% ---- pg-026 （书页 14） ----
一个三维散布图通常能用来揭示群组结构.

\mvbold{例 1.7} （三维情形下群组结构的搜寻）

参见例 1.6, 考虑雄性蜥蜴和雌性蜥蜴在三维数据空间中是否占据空间的不同部分, 是很有趣的. \hyperref[tab:1.3]{表 1.3} 中对应的数据中, 蜥蜴的性别(用行表示)分别为

\[
\begin{array}{cccccccccccc}
f & m & f & f & m & f & m & f & m & f & m & m \\
m & m & m & f & m & m & m & f & f & m & f & f
\end{array}
\]

\hyperref[fig:1.8]{图 1.8} 重复了原来的数据散布图, 只不过这里用空心点表示雌性, 用实心点表示雄性. 显然, 雄性的体型一般都比雌性的要大.

\begin{mvfigure}
\mvinc{pg-026}{0.26909,0.28730,0.73657,0.45937}
\phantomsection\label{fig:1.8}\mvcap{图 1.8\quad 雄性和雌性蜥蜴的三维散布图}
\end{mvfigure}

$n$ 维空间中 $p$ 个点 \quad $p$ 个变量的 $n$ 个观测结果也可看做是 $n$ 维空间中的 $p$ 个点. $\bm{X}$ 的每一列确定其中的一个点. 由第 $i$ 个变量的所有 $n$ 个观测值组成的第 $i$ 列确定了第 $i$ 个点

\[
\begin{bmatrix}
x_{1i} \\
x_{2i} \\
\vdots \\
x_{ni}
\end{bmatrix}
\]

在第 3 章中, 我们将介绍 $n$ 维空间中点的封闭性如何与相关变量之间联合的测量相联系.

\section{数据的展示及图表示}  % 原书 1.4
\phantomsection\label{sec:1.4}

强大的个人电脑与计算机工作站的迅速发展, 导致为数据分析和图解法而设计的复杂统计软件增加. 例如, 坐在书桌前就能通过灵巧的计算机生成的图片来观察多维数据的本质, 这已是可能的. 这些图片对于了解数据极有帮助, 并常常防止许多错误的出现以及随后的推理上的问题.

正如我们将在第 8 章和第 12 章看到的, 有许多方法企图用不多的维数表示 $p$ 维观测值, 使得每对观测值的初始距离(或相似性)几乎保持不变. 一般的, 如果多维观测值可以表现在二维空间中, 则离群值, 各种关系及可分辨的分组常常可以用眼睛辨别. 我们将讨论并说明几种在二维空间中显示多元数据的方法. 关于图解法的更多讨论的优秀初始资料见文献[11].
